% Template for PLoS
% Version 3.5 March 2018
%
% % % % % % % % % % % % % % % % % % % % % %
%
% -- IMPORTANT NOTE
%
% This template contains comments intended 
% to minimize problems and delays during our production 
% process. Please follow the template instructions
% whenever possible.
%
% % % % % % % % % % % % % % % % % % % % % % % 
%
% Once your paper is accepted for publication, 
% PLEASE REMOVE ALL TRACKED CHANGES in this file 
% and leave only the final text of your manuscript. 
% PLOS recommends the use of latexdiff to track changes during review, as this will help to maintain a clean tex file.
% Visit https://www.ctan.org/pkg/latexdiff?lang=en for info or contact us at latex@plos.org.
%
%
% There are no restrictions on package use within the LaTeX files except that 
% no packages listed in the template may be deleted.
%
% Please do not include colors or graphics in the text.
%
% The manuscript LaTeX source should be contained within a single file (do not use \input, \externaldocument, or similar commands).
%
% % % % % % % % % % % % % % % % % % % % % % %
%
% -- FIGURES AND TABLES
%
% Please include tables/figure captions directly after the paragraph where they are first cited in the text.
%
% DO NOT INCLUDE GRAPHICS IN YOUR MANUSCRIPT
% - Figures should be uploaded separately from your manuscript file. 
% - Figures generated using LaTeX should be extracted and removed from the PDF before submission. 
% - Figures containing multiple panels/subfigures must be combined into one image file before submission.
% For figure citations, please use "Fig" instead of "Figure".
% See http://journals.plos.org/plosone/s/figures for PLOS figure guidelines.
%
% Tables should be cell-based and may not contain:
% - spacing/line breaks within cells to alter layout or alignment
% - do not nest tabular environments (no tabular environments within tabular environments)
% - no graphics or colored text (cell background color/shading OK)
% See http://journals.plos.org/plosone/s/tables for table guidelines.
%
% For tables that exceed the width of the text column, use the adjustwidth environment as illustrated in the example table in text below.
%
% % % % % % % % % % % % % % % % % % % % % % % %
%
% -- EQUATIONS, MATH SYMBOLS, SUBSCRIPTS, AND SUPERSCRIPTS
%
% IMPORTANT
% Below are a few tips to help format your equations and other special characters according to our specifications. For more tips to help reduce the possibility of formatting errors during conversion, please see our LaTeX guidelines at http://journals.plos.org/plosone/s/latex
%
% For inline equations, please be sure to include all portions of an equation in the math environment.  For example, x$^2$ is incorrect; this should be formatted as $x^2$ (or $\mathrm{x}^2$ if the romanized font is desired).
%
% Do not include text that is not math in the math environment. For example, CO2 should be written as CO\textsubscript{2} instead of CO$_2$.
%
% Please add line breaks to long display equations when possible in order to fit size of the column. 
%
% For inline equations, please do not include punctuation (commas, etc) within the math environment unless this is part of the equation.
%
% When adding superscript or subscripts outside of brackets/braces, please group using {}.  For example, change "[U(D,E,\gamma)]^2" to "{[U(D,E,\gamma)]}^2". 
%
% Do not use \cal for caligraphic font.  Instead, use \mathcal{}
%
% % % % % % % % % % % % % % % % % % % % % % % % 
%
% Please contact latex@plos.org with any questions.
%
% % % % % % % % % % % % % % % % % % % % % % % %

\documentclass[10pt,letterpaper]{article}
\usepackage[top=0.85in,left=2.75in,footskip=0.75in]{geometry}

% amsmath and amssymb packages, useful for mathematical formulas and symbols
\usepackage{amsmath,amssymb}

% Use adjustwidth environment to exceed column width (see example table in text)
\usepackage{changepage}

% Use Unicode characters when possible
\usepackage[utf8x]{inputenc}

% textcomp package and marvosym package for additional characters
\usepackage{textcomp,marvosym}

% cite package, to clean up citations in the main text. Do not remove.
\usepackage{cite}

% Use nameref to cite supporting information files (see Supporting Information section for more info)
\usepackage{nameref,hyperref}

% line numbers
\usepackage[right]{lineno}

% ligatures disabled
\usepackage{microtype}
\DisableLigatures[f]{encoding = *, family = * }

% color can be used to apply background shading to table cells only
\usepackage[table]{xcolor}

% array package and thick rules for tables
\usepackage{array}

% create "+" rule type for thick vertical lines
\newcolumntype{+}{!{\vrule width 2pt}}

% create \thickcline for thick horizontal lines of variable length
\newlength\savedwidth
\newcommand\thickcline[1]{%
  \noalign{\global\savedwidth\arrayrulewidth\global\arrayrulewidth 2pt}%
  \cline{#1}%
  \noalign{\vskip\arrayrulewidth}%
  \noalign{\global\arrayrulewidth\savedwidth}%
}

% \thickhline command for thick horizontal lines that span the table
\newcommand\thickhline{\noalign{\global\savedwidth\arrayrulewidth\global\arrayrulewidth 2pt}%
\hline
\noalign{\global\arrayrulewidth\savedwidth}}


% Remove comment for double spacing
%\usepackage{setspace} 
%\doublespacing

% Text layout
\raggedright
\setlength{\parindent}{0.5cm}
\textwidth 5.25in 
\textheight 8.75in

% Bold the 'Figure #' in the caption and separate it from the title/caption with a period
% Captions will be left justified
\usepackage[aboveskip=1pt,labelfont=bf,labelsep=period,justification=raggedright,singlelinecheck=off]{caption}
\renewcommand{\figurename}{Fig}

% Use the PLoS provided BiBTeX style
\bibliographystyle{plos2015}

% Remove brackets from numbering in List of References
\makeatletter
\renewcommand{\@biblabel}[1]{\quad#1.}
\makeatother



% Header and Footer with logo
\usepackage{lastpage,fancyhdr,graphicx}
\usepackage{epstopdf}
%\pagestyle{myheadings}
\pagestyle{fancy}
\fancyhf{}
%\setlength{\headheight}{27.023pt}
%\lhead{\includegraphics[width=2.0in]{PLOS-submission.eps}}
\rfoot{\thepage/\pageref{LastPage}}
\renewcommand{\headrulewidth}{0pt}
\renewcommand{\footrule}{\hrule height 2pt \vspace{2mm}}
\fancyheadoffset[L]{2.25in}
\fancyfootoffset[L]{2.25in}
\lfoot{\today}

%% Include all macros below

\newcommand{\lorem}{{\bf LOREM}}
\newcommand{\ipsum}{{\bf IPSUM}}

%% END MACROS SECTION


\begin{document}
\vspace*{0.2in}

% Title must be 250 characters or less.
\begin{flushleft}
{\Large
\textbf\newline{Direct and indirect impact of fog on the hydrological balance of a subalpine forest} 
}
\newline
% Insert author names, affiliations and corresponding author email (do not include titles, positions, or degrees).
\\
Glenda García-Santos\textsuperscript{1},
Nikolaus Obojes\textsuperscript{2},
Polina Lemenkova\textsuperscript{3,4*},
Leonardo Montagnani\textsuperscript{3}
\\
\bigskip
\textbf{1} Department of Geography and Regional Studies, University of Klagenfurt, Klagenfurt, Carinthia, Austria
\\
\textbf{2} Institute for Alpine Environment, Eurac Research, Viale Druso 1, Bozen, Trentino-Alto Adige / South Tyrol, Italy
\\
\textbf{3} Department of Biological, Geological and Environmental Sciences, Alma Mater Studiorum -- Università di Bologna, Bologna, Emilia-Romagna, Italy
\\
\textbf{4} Faculty of Agricultural, Environmental and Food Sciences, Free University of Bozen-Bolzano, Piazza Università 1, 39100 Bolzano, Trentino-Alto Adige / South Tyrol, Italy
\\
\bigskip

% Insert additional author notes using the symbols described below. Insert symbol callouts after author names as necessary.
% 
% Remove or comment out the author notes below if they aren't used.
%

% Current address notes
%\textcurrency Current Address: Dept/Program/Center, Institution Name, City, State, Country % change symbol to "\textcurrency a" if more than one current address note
% \textcurrency b Insert second current address 
% \textcurrency c Insert third current address

% Use the asterisk to denote corresponding authorship and provide email address in note below.
* polina.lemenkova@unibz.it

\end{flushleft}
% Please keep the abstract below 300 words
\section*{Abstract}
Subalpine forests in the Alps are fragile ecosystems with high importance for human water resources and the local and mesoscale climate. While previous studies have measured different components of the water balance, little is known about the frequency and influence of fog and the role of forest age in the water balance. We conducted a comprehensive study of a subalpine coniferous forest at the site of Renon, in the Italian Alps, where a dense, old-growth forest section, intertwined with young patches, was monitored using the eddy covariance (EC) technique, tree transpiration sensors, phenocam images, throughfall and stemflow gauges, water discharge measurements, soil moisture sensors and epiphytes quantification. The interception was found to play a dominant role in the precipitation and evapotranspiration partitioning, especially in the older stand, where it was likely enhanced by the presence of lichens. Tree transpiration was surprisingly lower in the young stand and the evapotranspiration of soil and understory contributed considerably to the water balance at both stands. Fog caused additional throughfall in mixed fog and rain precipitation events. Discharge and the change of soil water content played a minor role in the yearly hydrological balance which was almost closed. This study revealed that precipitation interception and evapotranspiration partitioning change with forest age, and that fog plays a considerable role in the water balance of temperate, coniferous mountain forests, though it appears to be less frequent than in tropical and subtropical could forests.

% Please keep the Author Summary between 150 and 200 words
% Use first person. PLOS ONE authors please skip this step. 
% Author Summary not valid for PLOS ONE submissions.   
%\section*{Author summary}
%Lorem ipsum dolor sit amet, consectetur adipiscing elit. Curabitur eget porta erat. Morbi consectetur est vel gravida pretium. Suspendisse ut dui eu ante cursus gravida non sed sem. Nullam sapien tellus, commodo id velit id, eleifend volutpat quam. Phasellus mauris velit, dapibus finibus elementum vel, pulvinar non tellus. Nunc pellentesque pretium diam, quis maximus dolor faucibus id. Nunc convallis sodales ante, ut ullamcorper est egestas vitae. Nam sit amet enim ultrices, ultrices elit pulvinar, volutpat risus.

\linenumbers

% Use "Eq" instead of "Equation" for equation citations.
\section*{Introduction}

\subsection{General background}
The Alps represent an essential source of water availability in Europe, a role expected to become even more crucial under climate change conditions. Vegetation interacts profoundly with the water cycle, influencing both water availability for human needs and the distribution between sensible heat and latent heat, therefore interacting with the radiative forcing feedback mechanism. Mountain regions are of essential importance for water provision with an impact extending far wider than their actual range (Messerli et al., 2004; Viviroli et al., 2007; Beniston et al., 2011).  In the Alps and other Central European mountains, the subalpine elevation belt is mostly covered by forests dominated by conifers (Bohn et al., 2000). The eastern part of the Alps is particularly humid and is dominated by spruce forests. In South Tyrol, spruce forests cover 178,000 hectares, representing 52.8\% of all the forests in the region, which is covered by forests on 45.7\% of its surface (IFNC, 2022). However, they are also strongly affected by climate change with past and projected warming exceeding the global average (Beniston, 2005; Gobiet et al., 2014). 

Despite the importance of the interaction between forests and regional hydrology, few studies have successfully determined the various components of the water cycle and evapotranspiration partitioning at the basin level on forested terrain. In particular, given the high cloud cover characterising the Alpine mountains (XXX), we ask what the direct contribution of fog water might be to the hydrological cycle. Additionally, we explore whether fog might indirectly contribute to ecosystem water availability by promoting the growth of lichens, which can influence the water cycle due to their inherent water retention capacity.

\subsection{Water balance, discharge, and percolation}

Measurements of the water balance components are usually surrounded by considerable uncertainties (García-Santos, 2007) and, in hydrological approaches, widely accepted simplifications. Experimental studies where the different components of the water cycle are directly measured are scarce. Frequently, one or more of the components of the balance are assessed indirectly, with a modelling approach (XXX). The same happens concerning flux partitioning (Nelson et al., 2020). In ecohydrology, the partitioning is resolved by separating evapotranspiration into evaporation from surfaces inside the forest and transpiration by the vegetation (Mastrotheodoros et al., 2020; Savenije, 2004). While tree transpiration and soil evaporation can be measured directly through sapflow sensors and for instance canopy chambers, evaporation from canopy interception can still only be estimated indirectly or modelled. Thus, there are still major uncertainties, for example, concerning the rate of wet canopy evaporation estimated from eddy covariance measurements or as the differences between precipitation, throughfall and stemflow, or those using the conventional Penman-Monteith-equation by a factor of two or more (van Dijk et al., 2015; Holwerda et al., 2012). 

\subsection{Contribution of fog to water input}

Fog can represent an additional source of water, not quantified by the commonly used pluviometers, to the overall water input. In subtropical ecosystems with high and constant fog frequency, like the cloud forests in La Gomera island, on foggy days, the daily water budget may represent up to 49\% of the annual recorded water input (García-Santos, 2007; García-Santos et al. 2009; García-Santos and Bruijnzeel, 2011). However, the relevance of fog for the water balance of temperate mountains such as the Alps is still largely unknown. Eugster suggested that the alpine contribution of fog is limited, besides the frequent occurrence of low-elevation clouds (XXX). A distinct feature of alpine fog, as compared to fog occurring on islands, and promoting the occurrence of fog vegetation, is that fog frequently occurs together with liquid precipitation, determining the ‘mixed precipitation’ events (XXX). 

\subsection{Fraction of water transpired by the vegetation}

Another topic relevant to the understanding of subalpine forest ecohydrology is the evapotranspiration (ET) partitioning between evaporation (E) and transpiration (T). Recently, several models have been developed to partition ET (Nelson et al., 2018 , Zhou et al., 2016) and these models have been tested across several climates, plant functional types and sites (Nelson et al., 2020). Alpine forests, however, have their specificity. First, considering that the soil is covered by high vegetation with dense foliage, we can assume that soil evaporation represents a minor source of E as compared to wet canopy E, contrary to what happens, for instance, in sparse broadleaved forests like woody savanna. Secondly, Mountain forests have a specific combination of temperature and precipitation. In some cases, the combination of high precipitation and low temperature is unique, exceeding the combination of T and P depicted by the Whittaker scheme (Testolin et al., 2020). In a global study by van Dijk et al. (2015) the Renon forest, well representing the typical subalpine forest structure, was found to be the site with the highest degree of water intercepted by the canopy, therefore promoting a large fraction of wet canopy E.

There any differences between old and young forest formations in the water interception capacity. Nevertheless, the role of epiphytes in the canopy water retention capacity remains the question in forest ecology. From a forest management perspective, it would be relevant to understand whether young or old forests, growing in this specific habitat conditions, have a higher water retention capacity. Young forests generally achieve quite early a similar leaf area as mature forests, therefore having a similar water-holding capacity. Trunks and branches are however obviously larger in the old forest sections, but spruce and pine cork have limited capacity to hold water. Subalpine forests, however, frequently show an abundant cover of epiphytes, bryophytes and particularly lichens. Previous studies assessed their species composition and their interaction with the climate. However, little is known about their water-holding capacity and hence the extent of their interaction with the water cycle. Due to their poikilohydric nature, they are also expected to behave differently in terms of transpiration, presumably losing water without high water deficit limitation as the trees do. If relevant, this behaviour would affect the ecosystem response to energy partitioning (Porada et al., 2018). The presence of lichens can therefore represent an indirect effect of fog occurrence since it is extremely effective in providing water to these organisms, which can make photosynthesis and develop as a direct function of their water status (XXX).

\subsection{Objectives and goals}

In the context of the interaction of forests and meteorological drivers in conditions of climatic change and, we are aiming a better understanding of the ecohydrology of a subalpine forested catchment. In this paper, we detail the following working questions:

\begin{enumerate}
    \item What is the average fraction of water intercepted by vegetation?
    \item What is the direct contribution of fog to water input?
    \item What are the drivers of water interception and discharge?
    \item What is the fraction of water transpired by the vegetation?
    \item What is the role of forest in the partitioning of precipitation between water intercepted by the canopy and the soil and water discharge ?
    \item Are there any differences between old and young forest formations in the water interception capacity? What is the role of epiphytes in this difference?
 \end{enumerate}
 
To answer these questions, we complemented a highly instrumented long-term monitoring site, belonging to the ICOS network, where precipitation, ET through eddy covariance technique and SWC variation at several depths are routinely performed with additional measurements, including water discharge, sap flow, throughfall, fog interception and water holding capacity of lichens.

\section{Materials and methods}
\subsection{Study area}

The experiment was conducted at the Renon site, South Tyrol, in the Italian Alps (1735 m a.s.l., 46°35′11′′N, 11°26′00′′E). The site, IT-Ren, provides data for several monitoring networks, including Fluxnet (Pastorello et al., 2020), ICOS (https://www.icos-cp.eu/), and Lter (https://www.lter-europe.net/).  The catchment has an area of 0.44 km2. At the site, it is present a 40 m tall tower where eddy covariance measurements are routinely performed giving the information on evapotranspiration at a half-hour time step. The water basin where the experiment was conducted was measured using the local digital elevation model using ArcGIS software (ESRI, Redlands, CA, USA).
The soil has developed above a layer of glacial till, with a depth of approximately 1 m, placed on top of a porphyry bedrock. The soil was classified as Haplic Podzol according to the FAO soil taxonomy (FAO, 1998) and on average consisted of 49\% sand, 39\% silt, and 12\% clay.

The tree layer (diameter at breast height (DBH) >5 cm) consisted of 85\% spruce [Picea abies (L) Karst.], 12\% Swiss stone pine (Pinus cembra L.), and 3\% European larch (Larix europea L.) trees. Scots pine (Pinus sylvestris L.) and European rowan (Sorbus aucuparia) individuals were also present sparsely. The dominant tree height was approximately 29 m. The understory consisted mainly of blueberry (Vaccinium myrtillus L.). Intervening grasslands were dominated by wavy hair-grass [Deschampsia flexuosa (L.) Trin].

The forest is of natural origin and managed for wood production. The traditional harvest method creates small gaps, approximately 50 m wide, and involves the thinning of surrounding trees. The result is a heterogeneous vegetation structure, with groups that are almost even-aged forming an unevenaged structure at a larger scale. A large group of dominant spruce trees with an age of approximately 200 years, and a second group of young, about 30 years-old trees were present at the study site. In both stands, parts of the living crown frequently reached the ground, irrespective of age. The diameters and heights of the trees in the research area have been measured every ten years since 1990, while the diameter of a subset of trees is measured annually by manual dendrometers (UMS, München, Germany). The most recent inventory (tree height, size, and position) was performed in 2020 (Badraghi et al., 2021) with the TruPulse sensor (TruPulse 360 B laser range-finder; Laser Tech, Centennial CO, USA). In addition, tree size was assessed during the summer of 2020 using a laser technique involving a UAV-mounted laser scanner flight (Fig. 1 Supplementary information). Local spruce and Swiss stone pine were characterised by an almost columnar shape (Marcolla et al., 2005). This shape and the high leaf area index (LAI) (average 5.1 as measured by hemispherical images), created peculiar microclimatic conditions within the crown.

\subsection{Measurements of water balance at the catchment level }

The water balance at the catchment level can be quantified by the following equation: 

Where ET is evapotranspiration, P is Precipitation, C is capillarity, D is drainage, SW is soil water content and R is runoff.
Here, we considered and quantified all these components except C and D, including the fog, which is frequently neglected in water balance studies.

\paragraph{ET}
The evapotranspiration was assessed by the eddy covariance technique. The measuring system was composed of a 3D ultrasonic anemometer (Gill HS-50, Lymington, UK) and an enclosed-path infrared (IR) gas analyser (Li 7200 LiCor, Lincoln, NE, USA). The instruments were placed on a tower at 33.7 m above ground level.

Air samples were taken through an insulated steel tube of 4 mm in internal diameter and 0.75 m in length at 0.15 m from the anemometer. The airflow rate was set to 12 L min− 1, as provided by the flow module (7200–101, Li-Cor). All the setups followed the ICOS methodology (Rebmann et al., 2018). Raw values of CO2 and H2O concentration and 3D wind speed components were measured at 20 Hz and the resulting fluxes were computed and logged by a personal computer every 30 min. and elaborated with Eddypro software (version 6.2.1, LiCor). 

Based on the energy balance ratio (EBR) values computed monthly (Wohlfahrt et al., 2009), we estimated the systematic errors in the eddy covariance measurements for LE+H at the Renon site to be 17.6\%. Given a Bowen ratio of approximately 1 in the summer, this error can be assumed to be half of this value in LE (8.8\%; Table 1) and, therefore also in ET (Table 1). The random error of EC and the other components of the water balance are shown in Figure 8.

\paragraph{P}

Precipitation was assessed by a weighing pluviometer Geonor T200x. The pluviometer was installed on a steel basement placed in a forest gap. The data were collected on a CR3000. In parallel, a tipping-bucket pluviometer (Hobo, Onset, MA, USA) was also installed. The data were collected on CR6 Campbell Scientific datalogger. All the data were sampled every 20 seconds and averaged at 30 30-minute time steps.

\subsection{Instruments}

The role and contribution of fog were estimated by its impact on throughfall, following the method in Hutley et al. (1997). We first determined a linear regression equation between throughfall and precipitation rates for mixed precipitation and rain-only events for both stands. These gave a predicted throughfall for a given precipitation event with rain only and with fog (mixed precipitation). Fog contribution to throughfall for mixed precipitation days was then estimated as the difference between measured throughfall and the contribution of rain to throughfall calculated as the product of precipitation and the slope of the throughfall to precipitation equation from rain-only events. 

\subsection{Runoff}

Water discharge at the catchment scale was measured using a combination of a water stage sensor and flow velocity measurements performed with the salt dilution method. The water stage sensor was placed at the lowest spot of the catchment, just above an artificial water basin (46°35′00′′N, 11°26′02′′E, 1,675 m a.s.l.) and continuously measured the height of the water table in the stream (S in cm). Discharge (DC in l s-1) was calculated using the following equation:  

where a and b are empirical parameters according to the salt method (Day, 1977). 

Continuous water depth measurements were conducted in the creek with an OTT PLS 500 (OTT HydroMet Corp, VA, USA). Stream flow velocity was assessed monthly and in two more instances during 2019 after heavy precipitation events. 

\subsection{Sensors}

Tree transpiration was estimated by measuring the sap flow of five spruce trees with a DBH ranging from 23 cm to 57 cm. The sap flow measurements of up to 10 trees going back to 2016 showed a representative behaviour of these trees for their size classes. To minimise errors due to incoming shortwave radiation, one sensor of the tissue heat balance sensor (Table 1) was installed on the north side of the trees. The measuring system provided sap flow rates integrated for the whole sapwood depth per unit trunk circumference (kg h-1 cm-1). We logged the measured values at 10-min intervals and then scaled them to the tree level by multiplying them with stem circumference (minus the bark and phloem thickness measured during sensor installation) and integrated them to 30-min and daily sums per tree (L day-1) according to Kučera (2010). When a sensor was installed for more than a year, we checked whether wound reaction and/or an accumulation of resin led to a decrease in sap flow by comparing the 95th percentile (P95) of 30-min sap flow rates for 2019 and the previous years. For the two smallest trees, we found a considerable decrease in the P95 sap flow and consequently corrected sap flow for 2019 by multiplying it with the ratio of P95 [(year of sensor installation)/P95(2019)]. Minor data gaps caused by power outages or temporary sensor malfunction were filled by using a linear correlation of the respective tree’s sap flow with eddy covariance evapotranspiration (R² = 0.57 to 0.66 at a 30-min resolution). 

Tree transpiration was calculated by dividing sap flow by the projected crown area of the respective tree, thereby converting units from kg per tree to mm. The projected crown area was estimated from the mean crown radius measured in the four cardinal directions. The average tree transpiration of the two smaller trees (DBH 23 cm and 32 cm, respectively) was used to represent the young stand, while the average of the three larger trees (DBH 43 cm, 50 cm, and 57 cm, respectively) was used to represent the old stand. All the selected trees were dominant or codominant.

Precipitation below the canopy as throughfall was measured with sixteen manual rain gauges, arranged in two groups of eight in the two main forest formations, the 200-year-old section and the 30-year-old section. These pluviometers, with a 10-cm diameter orifice, were arranged in rows with a 5-m distance between each pluviometer and data were recorded almost on a weekly basis. Additionally, six tipping bucket pluviometers (with a Hobo data logger; Onset) arranged 5 m apart in two groups of three continuously recorded the below canopy precipitation at a 10-min resolution. To increase their collection representativeness, funnels with a diameter of 30 cm were placed above the orifice of each pluviometer (Figure 1).

\subsection{Impact of tree age and epiphytes}

The role of lichens in the water balance was assessed by measuring air temperature and humidity inside the crown and the water storage capacity of the lichens at two different tree heights (Table 1). The water storage capacity of the lichens was assessed in spring from a tree falling in the previous winter, with a tree height of 28 m and a DBH of 53 cm which is representative of the old stand. To estimate the lichen weight, the tree was divided into 3-m sections. In each section, all the branches were counted, and all lichens present above a single randomly selected branch were collected, together with the lichens growing on half of the main stem. In the laboratory, the lichens were first wetted until water saturation and then weighed to assess the fresh weight. They were then dried in an oven at 45 °C until a constant weight was achieved and then weighed (Sartorius Entris 2202, Göttingen, Germany) to assess their dry weight.

% Results and Discussion can be combined.
\section*{Results}
Nulla mi mi, venenatis sed ipsum varius, Table~\ref{table1} volutpat euismod diam. Proin rutrum vel massa non gravida. Quisque tempor sem et dignissim rutrum. Lorem ipsum dolor sit amet, consectetur adipiscing elit. Morbi at justo vitae nulla elementum commodo eu id massa. In vitae diam ac augue semper tincidunt eu ut eros. Fusce fringilla erat porttitor lectus cursus, vel sagittis arcu lobortis. Aliquam in enim semper, aliquam massa id, cursus neque. Praesent faucibus semper libero.

% Place tables after the first paragraph in which they are cited.
\begin{table}[!ht]
\begin{adjustwidth}{-2.25in}{0in} % Comment out/remove adjustwidth environment if table fits in text column.
\centering
\caption{
{\bf Table caption Nulla mi mi, venenatis sed ipsum varius, volutpat euismod diam.}}
\begin{tabular}{|l+l|l|l|l|l|l|l|}
\hline
\multicolumn{4}{|l|}{\bf Heading1} & \multicolumn{4}{|l|}{\bf Heading2}\\ \thickhline
$cell1 row1$ & cell2 row 1 & cell3 row 1 & cell4 row 1 & cell5 row 1 & cell6 row 1 & cell7 row 1 & cell8 row 1\\ \hline
$cell1 row2$ & cell2 row 2 & cell3 row 2 & cell4 row 2 & cell5 row 2 & cell6 row 2 & cell7 row 2 & cell8 row 2\\ \hline
$cell1 row3$ & cell2 row 3 & cell3 row 3 & cell4 row 3 & cell5 row 3 & cell6 row 3 & cell7 row 3 & cell8 row 3\\ \hline
\end{tabular}
\begin{flushleft} Table notes Phasellus venenatis, tortor nec vestibulum mattis, massa tortor interdum felis, nec pellentesque metus tortor nec nisl. Ut ornare mauris tellus, vel dapibus arcu suscipit sed.
\end{flushleft}
\label{table1}
\end{adjustwidth}
\end{table}


%PLOS does not support heading levels beyond the 3rd (no 4th level headings).
\subsection*{\lorem\ and \ipsum\ nunc blandit a tortor}
\subsubsection*{3rd level heading} 
Maecenas convallis mauris sit amet sem ultrices gravida. Etiam eget sapien nibh. Sed ac ipsum eget enim egestas ullamcorper nec euismod ligula. Curabitur fringilla pulvinar lectus consectetur pellentesque. Quisque augue sem, tincidunt sit amet feugiat eget, ullamcorper sed velit. Sed non aliquet felis. Lorem ipsum dolor sit amet, consectetur adipiscing elit. Mauris commodo justo ac dui pretium imperdiet. Sed suscipit iaculis mi at feugiat. 

\begin{enumerate}
	\item{react}
	\item{diffuse free particles}
	\item{increment time by dt and go to 1}
\end{enumerate}


\subsection*{Sed ac quam id nisi malesuada congue}

Nulla mi mi, venenatis sed ipsum varius, volutpat euismod diam. Proin rutrum vel massa non gravida. Quisque tempor sem et dignissim rutrum. Lorem ipsum dolor sit amet, consectetur adipiscing elit. Morbi at justo vitae nulla elementum commodo eu id massa. In vitae diam ac augue semper tincidunt eu ut eros. Fusce fringilla erat porttitor lectus cursus, vel sagittis arcu lobortis. Aliquam in enim semper, aliquam massa id, cursus neque. Praesent faucibus semper libero.

\begin{itemize}
	\item First bulleted item.
	\item Second bulleted item.
	\item Third bulleted item.
\end{itemize}

\section*{Discussion}
Nulla mi mi, venenatis sed ipsum varius, Table~\ref{table1} volutpat euismod diam. Proin rutrum vel massa non gravida. Quisque tempor sem et dignissim rutrum. Lorem ipsum dolor sit amet, consectetur adipiscing elit. Morbi at justo vitae nulla elementum commodo eu id massa. In vitae diam ac augue semper tincidunt eu ut eros. Fusce fringilla erat porttitor lectus cursus, vel sagittis arcu lobortis. Aliquam in enim semper, aliquam massa id, cursus neque. Praesent faucibus semper libero~\cite{bib3}.

\section*{Conclusion}

CO\textsubscript{2} Maecenas convallis mauris sit amet sem ultrices gravida. Etiam eget sapien nibh. Sed ac ipsum eget enim egestas ullamcorper nec euismod ligula. Curabitur fringilla pulvinar lectus consectetur pellentesque. Quisque augue sem, tincidunt sit amet feugiat eget, ullamcorper sed velit. 

Sed non aliquet felis. Lorem ipsum dolor sit amet, consectetur adipiscing elit. Mauris commodo justo ac dui pretium imperdiet. Sed suscipit iaculis mi at feugiat. Ut neque ipsum, luctus id lacus ut, laoreet scelerisque urna. Phasellus venenatis, tortor nec vestibulum mattis, massa tortor interdum felis, nec pellentesque metus tortor nec nisl. Ut ornare mauris tellus, vel dapibus arcu suscipit sed. Nam condimentum sem eget mollis euismod. Nullam dui urna, gravida venenatis dui et, tincidunt sodales ex. Nunc est dui, sodales sed mauris nec, auctor sagittis leo. Aliquam tincidunt, ex in facilisis elementum, libero lectus luctus est, non vulputate nisl augue at dolor. For more information, see \nameref{S1_Appendix}.

\section*{Supporting information}

% Include only the SI item label in the paragraph heading. Use the \nameref{label} command to cite SI items in the text.
\paragraph*{S1 Fig.}
\label{S1_Fig}
{\bf Bold the title sentence.} Add descriptive text after the title of the item (optional).

\paragraph*{S2 Fig.}
\label{S2_Fig}
{\bf Lorem ipsum.} Maecenas convallis mauris sit amet sem ultrices gravida. Etiam eget sapien nibh. Sed ac ipsum eget enim egestas ullamcorper nec euismod ligula. Curabitur fringilla pulvinar lectus consectetur pellentesque.

\paragraph*{S1 File.}
\label{S1_File}
{\bf Lorem ipsum.}  Maecenas convallis mauris sit amet sem ultrices gravida. Etiam eget sapien nibh. Sed ac ipsum eget enim egestas ullamcorper nec euismod ligula. Curabitur fringilla pulvinar lectus consectetur pellentesque.

\paragraph*{S1 Video.}
\label{S1_Video}
{\bf Lorem ipsum.}  Maecenas convallis mauris sit amet sem ultrices gravida. Etiam eget sapien nibh. Sed ac ipsum eget enim egestas ullamcorper nec euismod ligula. Curabitur fringilla pulvinar lectus consectetur pellentesque.

\paragraph*{S1 Appendix.}
\label{S1_Appendix}
{\bf Lorem ipsum.} Maecenas convallis mauris sit amet sem ultrices gravida. Etiam eget sapien nibh. Sed ac ipsum eget enim egestas ullamcorper nec euismod ligula. Curabitur fringilla pulvinar lectus consectetur pellentesque.

\paragraph*{S1 Table.}
\label{S1_Table}
{\bf Lorem ipsum.} Maecenas convallis mauris sit amet sem ultrices gravida. Etiam eget sapien nibh. Sed ac ipsum eget enim egestas ullamcorper nec euismod ligula. Curabitur fringilla pulvinar lectus consectetur pellentesque.

\section*{Acknowledgments}
Cras egestas velit mauris, eu mollis turpis pellentesque sit amet. Interdum et malesuada fames ac ante ipsum primis in faucibus. Nam id pretium nisi. Sed ac quam id nisi malesuada congue. Sed interdum aliquet augue, at pellentesque quam rhoncus vitae.

\nolinenumbers

% Either type in your references using
% \begin{thebibliography}{}
% \bibitem{}
% Text
% \end{thebibliography}
%
% or
%
% Compile your BiBTeX database using our plos2015.bst
% style file and paste the contents of your .bbl file
% here. See http://journals.plos.org/plosone/s/latex for 
% step-by-step instructions.
% 
\begin{thebibliography}{10}

\bibitem{bib1}
Conant GC, Wolfe KH.
\newblock {{T}urning a hobby into a job: how duplicated genes find new
  functions}.
\newblock Nat Rev Genet. 2008 Dec;9(12):938--950.

\bibitem{bib2}
Ohno S.
\newblock Evolution by gene duplication.
\newblock London: George Alien \& Unwin Ltd. Berlin, Heidelberg and New York:
  Springer-Verlag.; 1970.

\bibitem{bib3}
Magwire MM, Bayer F, Webster CL, Cao C, Jiggins FM.
\newblock {{S}uccessive increases in the resistance of {D}rosophila to viral
  infection through a transposon insertion followed by a {D}uplication}.
\newblock PLoS Genet. 2011 Oct;7(10):e1002337.

\end{thebibliography}



\end{document}

